\section{Conclusion}
We have presented a stateful NRX for 5G NR PUSCH with (i)~a pilot memory that carries LS channel estimates across slots, without labels and without decoder feedback, combined with a channel filter whose weights are fitted on the memory and not trained; (ii)~an evaluation in which, with two CGNN iterations, it clearly outperforms the two-iteration NRX retrained on the same data, outperforms the three-iteration NRX, performs close to the eight-iteration NRX, with and without power control, and outperforms it on TDL and urban macrocell channels; (iii)~a latency measurement by which it remains real-time capable, unlike the three-iteration NRX; and (iv)~two network decisions based on its evidence, of which falling back to a conventional receiver helps for few TBs, while with power control, the network can request most retransmissions before decoding finishes. Future work includes simulating the requested retransmissions, a comparison with classical channel estimation over several slots, and results on measured data and over multiple training seeds.
